\begin{center}
    {\large \textbf{Giorno 0-1} \textbar\ ¿Dónde está Ramón? \textbar\ 8-9 Agosto 2024 \textbar\ \textbf{Jakarta}, Java}
\end{center}

\noindent\rule{\textwidth}{0.4pt}
\begin{figure}[h!] % Portrait images below
    \centering
    % First portrait image (on the left)
    \begin{minipage}[h]{0.48\textwidth} % Set width equal to half the page
        \centering
        \includegraphics[width=\textwidth]{images/flags/Screenshot Jakarta.png}
    \end{minipage}
    \hfill
    % Text
    \begin{minipage}[h]{0.48\textwidth}
        \raggedright
        \textbf{Superficie}: 661,5 km² \\
        \vspace{0.3cm}
        \textbf{Popolazione}: 10,68 milioni \\
        \vspace{0.3cm}
        \textit{Jakarta, l'imponente capitale dell'Indonesia, sorge sulla costa nord-occidentale dell'isola di Java. Architettura, lingua e cucina sono state influenzate dalle culture che hanno vissuto nella città: giavanese, malese, cinese, araba, indiana ed europea.}
    \end{minipage}
\end{figure}

\landscapeLargeMargin{images/day1/PXL_20240808_112752327.MP.jpg}

\landscapeLargeMargin{images/day1/PXL_20240808_112728709.jpg}

\landscapeLargeMargin{images/day1/PXL_20240808_200722123.jpg}

\portraitHalfPage{images/day1/IMG-20240808-WA0013.jpg}

\landscapeLargeMargin{images/day1/PXL_20240809_135715757.jpg}

\noindent 8-9 Agosto 2024 \newline

In una soleggiata e afosa Torino ha inizio la seconda parte della nostra luna di miele. Scartata come meta a dicembre poiché fuori stagione, l'Indonesia è la nostra destinazione odierna. Il viaggio previsto per oggi è infinito, ma andiamo per punti.

Puntuali a metà mattina, i genitori di Elisa ci vengono a prendere a casa e ci accompagnano alla fermata del bus che ci porterà a Malpensa. O Aeroporto Silvio Berlusconi, per gli amici. Appena prendiamo posto, Elisa tira fuori dallo zaino le sue cremine che, con tanta pazienza — di entrambi — ha travasato in contenitori monoporzione nei giorni precedenti. Durante il viaggio attiviamo la eSIM indonesiana per poter rimanere connessi subito all'arrivo. Quest'anno partiamo preparati.

Giunti in aeroporto imbarchiamo velocemente i bagagli, superiamo i gate e ci prepariamo ad attendere l'ora della partenza in compagnia dei miei collaudati panini mozzarella e prosciutto da mezzo chilo l'uno.

È la prima volta che voliamo con Saudia Airlines e i comfort dell'aereo sono degni di nota: i sedili economy particolarmente comodi, la salvietta calda consegnata prima della partenza tramite una pinzetta sterile e la costante presenza della crew, che ininterrottamente provvede a distribuire cibo e bevande. Come note negative: la qualità del cibo e gli altoparlanti che trasmettono messaggi di ringraziamento ad Allah come benvenuto.

Fatte queste premesse, il primo volo è abbastanza faticoso. Eli dorme un po', io passo il tempo a leggere, ma entrambi patiamo il non poterci muovere in pieno giorno.

Atterriamo la sera a Jeddah, in Arabia Saudita. L'aeroporto sembra una piazza araba allestita con bancarelle al posto dei classici negozi. Nonostante la struttura avveniristica, riconosciamo ben pochi tratti occidentali, fatta eccezione per il McDonald's dove ci fermiamo per la cena, come ogni nostro spostamento che si rispetti.

Passiamo le poche ore di scalo camminando per l'aeroporto, per distendere le gambe in vista delle dieci ore che trascorreremo seduti sul prossimo volo. Mentre girovaghiamo, la nostra attenzione è attratta dalle innumerevoli aree di preghiera distribuite lungo tutto il perimetro dell'aeroporto e divise per uomini e donne.

Il nuovo aereo è esattamente come il primo, ma non facciamo in tempo a sentire i messaggi di ringraziamento ad Allah che entrambi crolliamo addormentati, vinti dalla stanchezza. Io sono periodicamente svegliato dalle assistenti di volo che, cercando di fare meno rumore possibile, ci adagiano bottigliette d'acqua e spuntini sulla ribaltina. A un certo punto ci consegnano un kit con spazzolino, dentifricio monodose, tappi e mascherina per gli occhi, e proprio quest'ultima mi permette di crollare in un lungo sonno ristoratore. Lungo per i miei standard, ovviamente, ma sufficiente a farmi recuperare un po' di forze.

A ormai poche ore dall'arrivo, sveglio Elisa: sta arrivando la colazione e so che non mi perdonerebbe mai di averla persa senza averla avvisata. Non mi perdona neanche di averla svegliata, ma al corso prematrimoniale ero stato avvisato su come funziona.

La scelta della colazione è tra riso con pollo, noodles con pollo e uova strapazzate. Chiediamo entrambi le uova, ma quando apriamo il contenitore troviamo i noodles con pollo. Piccanti, ovviamente.

Questo pasto per Elisa è traumatico. Il pollo piccante a colazione è decisamente troppo e lo assaggia appena. Punta allora tutto sul dolce e decide di prepararsi un panino con burro e marmellata di fragole. Fa un lavoro da certosino: stende un fazzoletto come tovaglietta, taglia minuziosamente la pagnotta col coltello di plastica, spalma prima il burro e poi la marmellata facendo attenzione a non uscire dai bordi del pane. Io provo a imitarla, ma il pane è talmente molle che quando tento di tagliarlo mi si frantuma in mano e finisco col mangiare la marmellata col cucchiaino. Nonostante riesca nella mastodontica impresa, il panino fa schifo e lo abbandona a metà. Prova con lo yogurt al mango, ma le basta assaggiare il mio per rinunciare: non aprirà neanche il suo vasetto.

Finalmente trova qualcosa di commestibile: la macedonia di frutta. Apre la monoporzione e fa per addentare il primo pezzo di melone, quando la passeggera spagnola alla sua sinistra — che finora era rimasta girata a parlare con la sua compagna di viaggio — si volta verso di lei, tossisce nella macedonia e ritorna a parlare con la vicina come se niente fosse. Elisa rimane col boccone di melone in mano e lo sguardo perso nel vuoto di chi è ormai morto dentro, mentre a me viene un attacco di ridarola che si ferma solo dopo vari minuti, quando ormai non riesco quasi più a respirare.

D'altro canto la mia colazione va un po' meglio. La fame vince sul gusto e inghiottisco tutto senza sentirne il sapore. Quando lo stomaco è pieno comincio a rendermi conto di quello che sto mangiando, ma per mia sfortuna questo accade mentre sto assaggiando lo yogurt al mango. Sembra di mangiare una gelatina di mango marcio: sarà l'unica cosa che avanzerò.

Assieme a noi sul volo si trova un gruppo di indonesiani in viaggio in qualche tour organizzato, vestiti tutti con gli stessi accessori rossi e bianchi: valigie, borse, sciarpe, persino una sedia a rotelle. Alla mia destra, un'anziana signora alta la metà di me trascorre tutto il viaggio tossendo a più non posso, e ogni volta che mi giro verso di lei intuisco che uno di noi due non arriverà a fine giornata. Dal momento che sto scrivendo questo diario, deduco che sia toccato a lei, e sono certo che la vecchietta mi mancherà.

Appena atterrati dobbiamo fare il Visa on Arrival, il visto per entrare in Indonesia, e ce la caviamo in pochi minuti. Andiamo poi di corsa a prendere i bagagli, perché vogliamo essere in hotel velocemente in modo da poter vedere Jakarta un paio d'ore prima di andare a riposare. Colpo di fortuna: i nostri due zaini sono sul nastro uno a fianco all'altro e li ritiriamo velocemente.

All'uscita prenotiamo subito un Grab, l'alternativa indonesiana di Uber, che si presenta puntuale in una decina di minuti e ci carica assieme ai bagagli per portarci in hotel. Davanti alla zona arrivi dell'aeroporto di Jakarta ci sono diverse stazioni di attesa per i servizi di noleggio e possiamo attendere il nostro passaggio su divanetti in pelle. Sono le sei di sera ma qui fa già buio.

Sul nostro Grab attraversiamo una parte di città periferica e ai lati dell'autostrada vedo un'immensa baraccopoli su più livelli che ricorda un po' le favelas di Rio de Janeiro. Quando arrivo in una nuova città la sera mi prende sempre un senso di sconforto. Era successo in Messico quando avevo guidato da Cancún fino alla nostra tenda al Pepem, col pensiero costante di essere fermato dalla polizia o di essere coinvolto in qualche incidente — e non aveva aiutato Elisa, che mi informava della possibilità di incontrare un giaguaro nella giungla.

Succede lo stesso anche qui, e mentre osservo questo scorcio povero di Jakarta passare attraverso il finestrino immagino come sarebbe avventurarsi per quelle strade e che genere di incontri si potrebbero fare. Sicuramente leggo troppi thriller, il che, unito alla stanchezza, fa sì che mi immagini legato su una sedia di metallo, nudo e sporco di sangue, in uno scantinato della baraccopoli a gridare "Non lo so dov'è Ramon!". Che poi neanche so chi sia Ramon. Forse è meglio guardare i grattacieli.

Il Mercure Jakarta Batavia, l'hotel per la notte, non ha niente a che vedere con la Jakarta periferica. È un'enorme struttura con tratti occidentali misto asiatici, con parcheggiatori e uscieri che ci accolgono cordiali mentre ci fanno passare sotto un metal detector all'ingresso, ignorando però tutti i bip. Facciamo il check-in, saliamo in una camera enorme — pagata pochissimo — con doccia di vetro in stanza, e ci prepariamo velocemente per uscire a fare un giro.

Tutto sembra essere andato per il verso giusto finché sento Elisa imprecare dopo aver trovato un lucchetto attaccato sul suo zaino. Inizialmente credo che lo abbiamo chiuso involontariamente e mi sembra un problema da poco, poi purtroppo Elisa mi mette di fronte alla realtà: abbiamo scambiato lo zaino con quello di qualcun altro. Quest'anno partiamo preparati, un par di palle.

Il modello dello zaino è esattamente uguale a quello di Elisa, ma manca la targhetta col nome e il contrassegno del check-in, oltre al già menzionato lucchetto di troppo. Seguono attimi di panico.

Per fortuna sappiamo che in queste situazioni la cosa da fare è reagire con lucidità, e cerchiamo di ricomporci subito. Dobbiamo tornare in aeroporto sperando che lo zaino sia ancora là. Per prima cosa però, sempre nell'ottica di mantenere la lucidità, scendiamo a cenare nella hall dell'albergo, perché a stomaco vuoto non si ragiona bene.

È una cena a buffet con prodotti locali e non va come previsto: Eli mangia solo riso in bianco, io provo un po' di tutto ma oltre al riso si salva solo il manzo. Provo una zuppetta di noodles, ma è piccante come il fuoco e non so quanto convenga bere l'acqua del distributore, perciò non riesco neanche a finirla. Mi gusto però un tè freddo che sarebbe anche stato buono se non fosse stato zuccheratissimo: non un'ottima combo con il manzo, insomma.

In pochi minuti siamo di nuovo su un Grab diretto verso l'aeroporto. La città è frenetica e molto trafficata, e quando il nostro loquace autista si ferma per un quarto d'ora a fare benzina imprechiamo silenziosamente entrambi. Alla pompa di fianco alla nostra osservo una fila di circa trenta motorini in attesa di fare rifornimento.

Durante il tragitto consideriamo i possibili scenari. L'altro passeggero potrebbe aver preso involontariamente lo zaino di Elisa e potrebbe stare per telefonarci per effettuare lo scambio. Oppure lo zaino potrebbe ancora essere in aeroporto nella zona ritiro bagagli. Potrebbe essere stato consegnato all'ufficio oggetti smarriti. O, ancora, l'altro passeggero — arrabbiato con noi — potrebbe averlo preso e buttato via. Quest'ultima è poco probabile, lo so, ma Elisa a volte è drastica.

Arrivati in aeroporto torniamo al Terminal 3 e spieghiamo la situazione agli addetti alla sicurezza. Sono tutti estremamente comprensivi e disponibili: ci controllano passaporto e biglietto e ci permettono di rientrare. Difficilmente mi sarei aspettato che questo fosse possibile in occidente, dal momento che non siamo passati attraverso alcun controllo con scanner o metal detector.

Giunti all'area bagagli, Eli si apposta davanti alla zona ritiro mentre io vado a cercare il reparto oggetti smarriti. Arrivo davanti al primo ufficio e — sorpresa! — lo zaino è lì! Corro a chiamare Elisa, che non sta più nella pelle. Appena lo trova, ci si avventa sopra e lo abbraccia, mentre io spiego all'addetto la situazione. Per fortuna avevamo scritto il nome sulla targhetta dello zaino, così dopo un rapido controllo dei documenti possiamo scambiare i due bagagli e uscire dall'aeroporto. Speriamo di non aver rovinato la vacanza a qualcuno.

Per il ritorno prenotiamo un altro Grab, che questa volta si fa attendere: il traffico davanti all'aeroporto è folle, ma gli addetti della Grab Pick-up Zone sono efficientissimi e si muovono tra le macchine come degli automi. Appena individuato il nostro mezzo, ci fiondiamo dentro. L'autista non parla inglese ed è una gioia, perché vogliamo solo riposare.

Calata la tensione del bagaglio, Elisa dorme secca mentre io mi trovo nuovamente a osservare la città passare dal finestrino. Al secondo giro quella sensazione di angoscia è già svanita e vedo la periferia sotto tutta un'altra luce. Ah, e mi ricordo anche chi è Ramon! L'autista messicano che ci ha accompagnato al Canyon del Sumidero e a Chiapa de Corzo! Non ti preoccupare Ramon, non ti tradirò.

Questi pensieri che si accavallano mi fanno realizzare che sono veramente stanco.

Arrivati in hotel, mentre Elisa si fa una doccia, sistemo tutti i dettagli per il viaggio del giorno seguente: prenoto il Grab per l'aeroporto, scarico le carte d'imbarco e confermo la nuova sistemazione a Lombok. Doccia veloce e andiamo a letto, collassando.

Welcome to Indonesia!

%\pagestyle{empty} % disable numbers


%\pagestyle{fancy} % enable numbers